%% sections/05-discussion.tex

\section{Discussion}
\label{sec:discussion}

\subsection{Styles Are Not Party-Archetype-Specific}

The most consequential result is the archetype-independence finding. All 7 styles fire
across all 7 campaigns, including Campaign 5 (professional motivations, no grief hooks),
Campaign 6 (duty-based institutional motivations), and Campaign 7 (minimal backstory,
no designed personal stakes). The styles originated in grief-themed campaigns (C1 and
C3 generated the largest instance counts), but their presence in C5--C7 confirms they
are not artifacts of the grief-themed design register.

This matters because it removes the primary alternative explanation for the styles:
that they are not design-enabled behavioral patterns but rather the predictable output
of campaigns designed around emotionally resonant content. If sheet-deep-reader only
fired in campaigns with detailed grief paragraphs, the style would not be a general
design principle; it would be a description of what grief-themed campaigns tend to
produce. The C5--C7 results rule this out.

The practical implication is that designers working in any motivational register ---
professional, institutional, relational, or grief-themed --- can apply the style
prescriptions. The prescriptions are not genre-specific; they are sheet-design-specific.

\subsection{The Promotion Threshold Is Conservative}

The 3-instance, 2-session threshold for promotion is more demanding than the rubric
amendment threshold (2 instances in a dimension). This conservatism was deliberate:
player styles are claimed to be \textit{ratified design patterns}, not provisional
observations, and the threshold was calibrated to exclude single-session anomalies.

In retrospect, the threshold appears well-calibrated. All 7 promoted styles accumulated
substantially more than 3 instances before ratification (the lowest-count style at
promotion, bilateral-arc-completion, had 5 instances). No style was promoted on the
threshold minimum alone. This suggests the threshold is functioning as a conservative
gate rather than as a certification mechanism: styles that are real accumulate well
beyond the minimum before anyone proposes them.

A potential objection is that the threshold is too conservative: genuinely useful design
patterns might be discarded for failing to reach 3 instances before the observation
window closes. The corpus has one candidate: the documentary-witness style was observed
in Campaign 4 but had only 1 instance before Session 7 of Campaign 4. Had the corpus
ended at Campaign 4, it might have been missed. The solution is a longer observation
window rather than a lower threshold: a 1-instance pattern that does not accumulate
across subsequent campaigns is not yet established as a design pattern.

\subsection{The Trigger-Type Shift Finding}

The RV11 confirmation --- that trigger type shifts from situational to relationship-driven
across Campaign 7's arc --- is a secondary finding with significant design implications.
It suggests that player styles are not static once established: the same style (say,
act-without-announcement) can be triggered by different types of events depending on
where the campaign is in its arc.

In Campaigns 1--3, trigger-type shifts of this kind were present but not analyzed
systematically. Campaign 7's explicit testing of RV11 provides the first systematic
evidence. The finding is consistent with a model in which PCs accumulate in-campaign
relational history across sessions and that history becomes progressively more available
as trigger material. Early in a campaign, PCs respond to scenes; late in a campaign,
they respond to each other.

This has adventure-design implications. A designer who wants act-without-announcement
instances in later sessions should design naming-event opportunities in earlier sessions
rather than scene-level triggers in the session where the behavioral change is wanted.
The trigger must precede the instance by at least one session.

\subsection{Limitations}

\paragraph{Single AI system.}
All 49 sessions were run by a single AI system playing PCs according to Playstyle
Heuristics. The styles are well-validated within this system, but generalization to
other AI systems or to human players requires further evidence. Different AI systems
may adhere to heuristics with different fidelity, which would affect instance rates.
Human players would produce additional instances not captured by the heuristics and
fewer instances that are purely heuristic-driven.

\paragraph{Corpus size.}
49 sessions across 7 campaigns is a substantial corpus for this type of analysis, but
it is not large by standards of quantitative game studies. The instance counts for
the lower-frequency styles (documentary-witness: 11 instances; intra-session
act-without-announcement: 8 instances) are not large enough to make strong statistical
claims. The findings should be treated as design hypotheses supported by qualitative
evidence rather than as statistically robust findings.

\paragraph{Single design team.}
All campaigns were designed by the same design pipeline, which has accumulated
conventions for PC sheet construction (dense grief paragraphs, named daily practices,
prior-adventure hooks) that may not be standard in other design contexts. Campaigns
designed by other teams with thinner PC sheets might produce lower instance counts,
particularly for sheet-deep-reader and craft-witness.

\paragraph{RV8 partial confirmation.}
The partial confirmation of RV8 (verbal-documentation-arc in a duty-trigger campaign)
suggests the style may have a design-dependency on oral-practice PC sheets that the
formal-reporting variant does not satisfy. This is the one area where the claim of
full archetype-independence is not yet established.
